% TOP_PROBLEM: {"id": 8700048, "title": "Conjecture 4.3 — Amoroso-David", "category_id": 3, "category": {"id": 3, "name": "graph_theory", "display_name": "Graph Theory", "description": "Problems involving graphs, networks, and their properties.", "slug": "graph-theory", "order_index": 3, "created_at": "2026-07-31T15:26:25.671Z"}, "status": "open", "problem_number": "AMR-086-0048", "set_id": 13}
% FIRST_POSTED: 2026-10-01
% ENTRY_KIND: statement_only
% UnsolvedMath ID 8700048; source catalogue code AMR-086-0048
% !TeX program = lualatex
% Definition and sources only; this file is not an LLM solution attempt.
\documentclass[11pt,a4paper]{article}
\usepackage[margin=25mm]{geometry}
\usepackage{fontspec}
\setmainfont{lmroman10-regular.otf}[BoldFont=lmroman10-bold.otf,
 ItalicFont=lmroman10-italic.otf,BoldItalicFont=lmroman10-bolditalic.otf]
\usepackage{amsmath,amssymb,mathtools}
\usepackage{unicode-math}
\setmathfont{latinmodern-math.otf}
\AtBeginDocument{\let\setminus\smallsetminus}
\usepackage{xurl,hyperref}
\hypersetup{colorlinks=true,urlcolor=blue}
\setcounter{secnumdepth}{0}
\setlength{\parindent}{0pt}
\setlength{\parskip}{5pt}
\setlength{\emergencystretch}{3em}
\begin{document}
\section*{Conjecture 4.3 — Amoroso-David}\label{8700048}
{\small UnsolvedMath ID 8700048 · AMR-086-0048 · Graph Theory · AMR Open Problem Lists}
\subsection{Definitions and mathematical statement}
Let $n\ge1$ and let $\alpha_1,\ldots,\alpha_n$ be nonzero algebraic complex numbers. They are multiplicatively independent if
\[
\alpha_1^{m_1}\cdots\alpha_n^{m_n}=1,\quad(m_1,\ldots,m_n)\in\mathbb Z^n
\]
forces $m_1=\cdots=m_n=0$. In particular, no entry is a root of unity. Put $D=[\mathbb Q(\alpha_1,\ldots,\alpha_n):\mathbb Q]$, the degree of their joint number field.

For an algebraic number $\alpha$ of degree $d$ with primitive minimal polynomial $a_d\prod_{j=1}^d(X-\alpha^{(j)})$ and $a_d>0$, define
\[
h(\alpha)=\frac1d\left(\log a_d+\sum_{j=1}^d
\log\max\{1,|\alpha^{(j)}|\}\right).
\]
This is the absolute logarithmic Weil height and fixes its normalization.

Does there exist, for each $n$, a constant $c(n)>0$ such that every multiplicatively independent $n$-tuple satisfies
\[
\prod_{i=1}^n h(\alpha_i)\ge\frac{c(n)}D?
\]
The constant may depend on the length of the tuple only. It cannot depend on the particular number field, the individual degrees, or the chosen algebraic numbers.

The denominator is the degree of the joint field, not the product of the individual degrees. The left side is a product, so an exceptionally small height must be compensated by the others in the precise way asserted. No integrality, positivity, or reality condition is imposed on the algebraic numbers. The case $n=1$ recovers the usual degree-dependent height form of Lehmer's problem.
\subsection{Short English statement}
Does a multiplicatively independent tuple of algebraic numbers satisfy the stated lower bound for its product of heights?
\subsection{Sources}
\begin{itemize}
\item[S1] ULAM.AI and the UnsolvedMath Contributors, supplied problems.json, record 8700048 (AMR-086-0048), AMR Open Problem Lists. Catalogue identity and source transcription; CC BY 4.0.
\url{https://huggingface.co/datasets/ulamai/UnsolvedMath}.
\item[S2] Waldschmidt - Open Diophantine Problems (2004), Conjecture 4.3. Original source indexed by AMR-086-0048.
\url{https://arxiv.org/abs/math/0312440}.
\end{itemize}
\end{document}
